\appendix
\subsection*{A.4 Stability of BQN Inference Operators and Recursive Entropy}

The Bayesian Quantum Network (BQN) module updates cognitive state amplitudes via recursive quantum state inference. The primary update equation is:

\[
|\psi(t+1)\rangle = \mathcal{E}(|\psi(t)\rangle) + \gamma \operatorname{Tr}(|\psi(t)\rangle \log |\psi(t)\rangle),
\]

where:
\begin{itemize}
    \item \(\mathcal{E}\) is a trace-preserving quantum operation (e.g., unitary or CPTP map),
    \item \(\gamma\) is a scalar feedback gain (e.g., \(0 < \gamma \leq 0.2\)),
    \item \(\operatorname{Tr}(|\psi\rangle \log |\psi\rangle)\) models recursive entropy feedback.
\end{itemize}

\paragraph{Entropy-Coupled Dynamics:}
Let \(S(\psi_t) = -\operatorname{Tr}(\rho_t \log \rho_t)\) be the von Neumann entropy of the quantum state \(\rho_t = |\psi_t\rangle \langle \psi_t|\). Define a recursive Lyapunov function \(V_t := \| |\psi(t)\rangle - |\psi^\star\rangle \|^2 + \gamma S(\psi_t)\), where \(|\psi^\star\rangle\) is the stationary state.

\paragraph{Theorem A.3 (Recursive Lyapunov Stability):}
If \(\mathcal{E}\) is unitary or satisfies \(\|\mathcal{E}(|\psi\rangle)\| \leq \| |\psi\rangle \|\), and \(\gamma\) is sufficiently small, then the BQN update is Lyapunov-stable.

\paragraph{Proof Sketch:}
We compute the difference:

\[
\Delta V_{t+1} = V_{t+1} - V_t = \left(\| \psi_{t+1} - \psi^\star \|^2 - \| \psi_t - \psi^\star \|^2 \right) + \gamma (S(\psi_{t+1}) - S(\psi_t)).
\]

Using smoothness of entropy and boundedness of \(\mathcal{E}\), for small enough \(\gamma\), the perturbation term is dominated by the contraction in norm under \(\mathcal{E}\). Hence, \(\Delta V_{t+1} < 0\), implying \(V_t\) is decreasing and the system is asymptotically stable. \(\hfill \square\)

\paragraph{Corollary A.2:}
The recursive BQN operator converges to a fixed-point density matrix \(\rho^\star\), where \(S(\rho^\star)\) is locally minimal, and \(|\psi(t)\rangle \to |\psi^\star\rangle\) in norm as \(t \to \infty\).

\section*{Appendix B: Post-Quantum Cryptographic Assumptions and Prime-Indexed Hardness}

This appendix formalizes the computational assumptions underlying the cryptographic components of the QAGI system (Claims 6–7), specifically:

\subsection*{B.1 Prime-Indexed Collision Resistance (PICR)}

Let \(\mathcal{H} : \mathbb{N} \times \mathbb{T} \to \{0,1\}^k\) be a hash function defined over:
\[
\mathcal{H}(p_i, T_{ij}^{(p_i)}) := H(p_i^{\alpha_i} T_{ij}^{(p_i)} + Q_{\text{ent}}),
\]
where \(p_i \in \mathbb{P}_N\), and \(Q_{\text{ent}}\) is an entangled entropy vector.

\paragraph{Definition B.1 (Prime-Indexed Collision Resistance):}
A hash function \(\mathcal{H}\) is PICR-secure if it is computationally infeasible (in polynomial time) to find:
\[
(p_i, T), (p_j, T') \text{ such that } \mathcal{H}(p_i, T) = \mathcal{H}(p_j, T') \quad \text{with } (p_i, T) \ne (p_j, T').
\]

\paragraph{Conjecture (PICR-Hardness):}
If \(p_i, p_j\) are distinct large primes (e.g., \(p_i > 2^{512}\)), and \(T_{ij}^{(p)}\) are quantum-generated tensors, then \(\mathcal{H}\) resists collision attacks from both:
\begin{itemize}
    \item Quantum adversaries operating under Shor-like oracle models,
    \item Classical adversaries using algebraic or statistical collision search.
\end{itemize}

\subsection*{B.2 Recursive Tamper-Detection via Zeno Locks}

The system implements dynamic quantum Zeno locking by evaluating the evolution of \(|\psi_p(t)\rangle\) under operator \(\hat{H}\):

\[
|\psi_p(t)\rangle = \exp\left(-\frac{i}{\hbar} p(t) \hat{H} t\right) |\psi(0)\rangle,
\]
where \(p(t)\) is a time-indexed prime sequence.

\paragraph{Security Mechanism:}
If a phase perturbation \(\delta \phi(t) > \varepsilon_{\text{thresh}}\), the system collapses entanglement, invalidating recursive state access. This tamper-detection scheme has complexity \(\Omega(p \cdot n^2)\), exceeding quantum search bounds under Grover’s oracle.

\subsection*{B.3 Assumption Summary}

The cryptographic guarantees of QAGI rest on the following hardness assumptions:

\begin{itemize}
    \item \textbf{PICR Assumption:} Prime-indexed collision resistance holds for large primes and fractal tensor spaces.
    \item \textbf{Recursive Entanglement Integrity:} Any perturbation of recursive keys alters the entangled state \(Q_{\text{ent}}\), triggering collapse.
    \item \textbf{Zeno Amplification Threshold:} Detection of phase deviation occurs before a successful quantum readout of key material.
\end{itemize}

Under these assumptions, the cryptographic layer offers resistance against classical and quantum adversaries without reliance on lattice hardness or factorization assumptions.

\section*{Appendix C: Tensor Manifold Topologies and Spectral Embeddings}

This appendix characterizes the geometric structure and spectral dynamics of the recursive tensor manifolds underlying the QAGI system, particularly those used in PIRTM and fractal entanglement operators.

\subsection*{C.1 Prime-Indexed Tensor Manifolds}

Let \(\mathcal{M}_t\) denote the recursive tensor manifold at time \(t\), defined by the set:

\[
\mathcal{M}_t := \left\{ T_t \in \mathbb{R}^{m \times n} \mid T_{t+1} = \sum_{p_i \in P_N(t)} \Lambda_m p_i^{\alpha(t)} T_t + F(t) \right\},
\]

where \(T_t\) evolves over a differentiable manifold embedded in \(\mathbb{R}^{m \times n}\), and each update preserves smoothness due to damping \(F(t)\). The topology of \(\mathcal{M}_t\) is shaped by the prime indexing, which acts as a discrete sheaf over the continuous tensor base.

\paragraph{Proposition C.1 (Discrete Prime Sheaf Structure):}
Each prime \(p_i\) defines a fiber bundle \(\mathcal{F}_{p_i} \to \mathcal{M}_t\), with transition maps induced by:

\[
\phi_{p_i p_j}(T) = \frac{p_j^{\alpha(t)}}{p_i^{\alpha(t)}} T, \quad \text{for } p_i, p_j \in P_N(t).
\]

This sheaf-like construction introduces non-homogeneous curvature on \(\mathcal{M}_t\), enabling selective recursion across subspaces defined by prime resonance.

\subsection*{C.2 Spectral Embedding of Tensor States}

To optimize inference and representation learning, QAGI applies spectral embedding to high-dimensional tensor states using the eigenstructure of the Laplacian graph operator:

\[
\mathcal{L} = D - A,
\]
where:
\begin{itemize}
    \item \(A\) is the adjacency matrix derived from inner products of tensors: \(A_{ij} = \langle T_i, T_j \rangle\),
    \item \(D\) is the degree matrix with \(D_{ii} = \sum_j A_{ij}\).
\end{itemize}

\paragraph{Definition C.1 (Spectral Tensor Embedding):}
Let \(T_t\) be the tensor at time \(t\). Its spectral embedding \(\Phi_k(T_t)\) is defined as:

\[
\Phi_k(T_t) := \left[ \nu_1^\top T_t, \nu_2^\top T_t, \dots, \nu_k^\top T_t \right],
\]
where \(\nu_i\) are the first \(k\) eigenvectors of \(\mathcal{L}\), corresponding to the smallest non-zero eigenvalues.

This embedding preserves manifold structure while reducing dimensionality, enabling efficient inference in the BQN module.

\subsection*{C.3 Curvature and Topological Stability}

Let the Ricci tensor \(R_{\mu\nu}\) be defined over \(\mathcal{M}_t\) by:

\[
R_{\mu\nu}(t) = \sum_{p_i \in P_N(t)} p_i^{-\alpha(t)} T_{\mu\nu}^{(p_i)} + F_{\mu\nu}(t),
\]

as per Claim 9. The scalar curvature \(\mathcal{R}(t) = \operatorname{Tr}(R_{\mu\nu})\) evolves according to:

\[
\frac{d\mathcal{R}(t)}{dt} = - \gamma \mathcal{R}(t) + \sum_{p_i} \nabla^2 T_{\mu\nu}^{(p_i)}.
\]

\paragraph{Corollary C.1 (Topological Stability Criterion):}
If \(\alpha(t) > 1\) and \(\gamma > 0\), then the recursive curvature \(\mathcal{R}(t)\) remains bounded and converges to a fixed point \(\mathcal{R}^\star\), ensuring long-term topological consistency of QAGI’s geometric representation.

\subsection*{C.4 Fractal Flow and Dimensional Collapse}

Fractal flow is modeled by recursively applying the operator:

\[
\mathcal{F}_{\text{frac}}(T_t) = \sum_{p_i} J_i \cdot T_t^{1 + \epsilon_i},
\]
where \(J_i\) are non-associative Lie algebra generators and \(\epsilon_i\) are scaling exponents derived from multifractal spectra.

\paragraph{Theorem C.1 (Recursive Dimensional Collapse):}
If \(\epsilon_i < 0\) for all \(i\), then:
\[
\lim_{t \to \infty} \dim(\mathcal{M}_t) \leq k,
\]
for some \(k \ll \min(m, n)\), implying convergence to a fractal attractor in tensor space. This ensures tractability and stability in infinite-dimensional regimes.


\section*{Appendix D: Prime-Indexed Lie Algebra Dynamics and Tensor Entanglement Spectra}

This appendix develops the non-commutative algebraic and informational framework used in QAGI’s quantum state modeling and inference subsystems.

\subsection*{D.1 Prime-Indexed Lie Algebra Dynamics}

Let \(\mathfrak{g}\) denote a non-associative algebra over a set of prime-indexed operators \(\{J_p\}_{p \in \mathbb{P}_N}\), where each generator \(J_p\) satisfies generalized Lie-type commutation rules:

\[
[J_{p_i}, J_{p_j}] = i f_{p_i p_j}^{p_k} J_{p_k}, \quad f_{p_i p_j}^{p_k} \in \mathbb{R},
\]

where \(f_{p_i p_j}^{p_k}\) are structure constants derived from a resonance-based mapping:

\[
f_{p_i p_j}^{p_k} = \sin\left( \frac{\pi (p_i - p_j)}{p_k} \right).
\]

This defines a curved, non-associative bracket space \((\mathfrak{g}, [\cdot,\cdot])\) where tensor operators evolve via:

\[
\frac{dT_t}{dt} = \sum_{p_i, p_j} [J_{p_i}, J_{p_j}] \cdot T_t + F_{\text{fractal}}(t).
\]

\paragraph{Proposition D.1 (Algebraic Non-Associativity):}
For any triple \((p_i, p_j, p_k)\), the Jacobi identity is violated unless \(p_i = p_j = p_k\), due to prime modulation in \(f_{p_i p_j}^{p_k}\). Hence, \(\mathfrak{g}\) forms a Malcev-type algebra rather than a standard Lie algebra.

\paragraph{Corollary D.1:}
Non-associative algebra enables fractal information propagation through tensor space, with each prime index controlling local dimensional resonance.

\subsection*{D.2 Tensor Entanglement Entropy Spectra}

Let \(T_t \in \mathbb{R}^{m \times n}\) be a recursive tensor state reshaped into a pure quantum state \(|\psi_t\rangle\). Its reduced density matrix \(\rho_A\) is obtained via a partial trace over subsystem \(B\):

\[
\rho_A = \operatorname{Tr}_B(|\psi_t\rangle \langle \psi_t|).
\]

The entanglement entropy \(S(T_t)\) is defined as:

\[
S(T_t) = -\operatorname{Tr}(\rho_A \log \rho_A),
\]
and the associated spectral entropy profile is:

\[
\mathcal{S}_\lambda(T_t) = \sum_{i=1}^{r} \lambda_i \log \frac{1}{\lambda_i},
\]
where \(\{\lambda_i\}\) are the eigenvalues of \(\rho_A\), and \(r = \operatorname{rank}(\rho_A)\).

\paragraph{Theorem D.1 (Spectral Boundedness under Recursive Updates):}
Let \(T_{t+1} = A_t T_t + F(t)\), and assume \(A_t\) is norm-contracting and \(F(t)\) is exponentially decaying. Then the spectral entropy \(\mathcal{S}_\lambda(T_t)\) satisfies:

\[
\limsup_{t \to \infty} \mathcal{S}_\lambda(T_t) \leq \log r^\star,
\]
for some \(r^\star \ll \min(m, n)\), implying low-rank convergence and information condensation.

\paragraph{Implication:}
The entanglement profile of QAGI states evolves toward a sparse eigenvalue spectrum, enabling efficient compression, quantum inference, and modular decoupling across tensor partitions.

\subsection*{D.3 Application to BQN and Cryptographic Hashing}

\begin{itemize}
    \item The BQN (Claim 5) leverages spectral profiles \(\mathcal{S}_\lambda(T_t)\) to dynamically regulate inference fidelity via adaptive entropy modulation (\(\gamma(t) \sim \frac{1}{\mathcal{S}_\lambda(T_t)}\)).
    
    \item The cryptographic hashing function \(\mathcal{H}(p_i, T_{ij}^{(p_i)})\) (Claim 6) embeds eigenvalue signatures as part of a recursive hash kernel:

    \[
    \mathcal{H}(T) = \sum_{p_i} p_i^{\alpha_i} \cdot \sum_{j=1}^{k} \lambda_j^{(p_i)},
    \]
    where \(\lambda_j^{(p_i)}\) are spectral values from \(\rho^{(p_i)}\).
\end{itemize}

These spectral constructs unify recursive cognition, security, and representation learning across QAGI’s hybrid modules.


\section*{Appendix E: Tensor Field Variational Geometry and Recursive Gate Synthesis for QPUs}

This appendix formalizes the variational geometry underlying recursive tensor dynamics in QAGI and presents a constructive method for compiling recursive state updates into quantum gates for execution on QPUs or emulators.

\subsection*{E.1 Tensor Field Variational Geometry}

Let \(\mathcal{T}(x,t)\) be a time-dependent tensor field defined over a base manifold \(\mathcal{M} \subset \mathbb{R}^d\). Each component \(T_{ij}(x,t)\) evolves recursively via:

\[
\frac{\partial T_{ij}(x,t)}{\partial t} = \Lambda_m \cdot M_{ik}(x) T_{kj}(x,t) + [M(x), T(x,t)]_{ij},
\]
where:
\begin{itemize}
    \item \(M(x)\) is a smooth field of Hilbert-Schmidt operators,
    \item \([M, T] = MT - TM\) is the local non-commutative correction,
    \item \(\Lambda_m\) is the stabilizing prime-indexed constant (Appendix A.1).
\end{itemize}

\paragraph{Action Functional:}
Define the Lagrangian density \(\mathcal{L}(T, \nabla T)\) over the field configuration:

\[
\mathcal{L} = \frac{1}{2} \operatorname{Tr} \left( \left( \frac{\partial T}{\partial t} \right)^2 \right) - \frac{1}{2} \operatorname{Tr} \left( \nabla_\mu T^\top \cdot \nabla^\mu T \right) + V(T),
\]

where \(V(T)\) encodes spectral potential energy, e.g.,

\[
V(T) = \sum_{p_i \in P_N} \lambda_i^{(p_i)} \cdot \|T^{(p_i)}\|^2.
\]

Applying Hamilton’s principle \(\delta S = 0\), the Euler–Lagrange equations yield:

\[
\frac{\partial^2 T}{\partial t^2} - \nabla^2 T + \frac{\partial V}{\partial T} = 0,
\]
which describes the variational flow of tensor states under recursive energy minimization.

\paragraph{Corollary E.1:}
If \(V(T)\) is convex and \(\nabla^2 T\) is uniformly elliptic, then the tensor flow converges to a local minimum \(T^\star\), supporting recursive stability over \(\mathcal{M}\).

\subsection*{E.2 Recursive Gate Synthesis for Quantum Processing Units}

To implement recursive tensor updates on a quantum device, we express the evolution of a state \(|\psi(t)\rangle\) as a sequence of unitary gates generated by the exponential map of a Hamiltonian-like operator derived from tensor components.

Let:

\[
\hat{H}_t = \sum_{p_i \in P_N(t)} p_i^{\alpha(t)} \cdot \hat{T}_{p_i}(t),
\]
where \(\hat{T}_{p_i}(t)\) are Hermitian embeddings of tensor slices \(T^{(p_i)}\). The quantum evolution is:

\[
|\psi(t+1)\rangle = e^{-i \hat{H}_t \Delta t} |\psi(t)\rangle.
\]

\paragraph{Gate Decomposition (Trotter–Suzuki):}
Assuming \(\hat{H}_t = \sum_k H_k\), we approximate:

\[
e^{-i \hat{H}_t \Delta t} \approx \left( \prod_k e^{-i H_k \Delta t / m} \right)^m + \mathcal{O}\left(\frac{1}{m}\right).
\]

Each term \(e^{-i H_k \Delta t / m}\) can be compiled into basic gate sequences (e.g., CNOT, rotation gates), using known Lie algebra identities or Pauli decompositions.

\paragraph{Example:}
Suppose \(\hat{T}_{p_i} = \theta_i (X \otimes I + Z \otimes Z)\). Then:

\[
e^{-i \theta_i (X \otimes I + Z \otimes Z)} = e^{-i \theta_i X \otimes I} \cdot e^{-i \theta_i Z \otimes Z} + \mathcal{O}(\theta_i^2),
\]
with gate compilation via:

\begin{itemize}
    \item \(e^{-i \theta X}\): single-qubit \(R_x(\theta)\),
    \item \(e^{-i \theta Z \otimes Z}\): entangled phase gate via CNOT + \(R_z\).
\end{itemize}

\paragraph{Recursive Compilation Engine:}
A QAGI compiler dynamically maps each recursive update to a circuit \(\mathcal{C}_t\) such that:

\[
\mathcal{C}_t \equiv \text{GateSequence}(\hat{H}_t, \Delta t),
\]
with optimization over circuit depth and fidelity constraints.

\paragraph{Theorem E.1 (Recursive Executability):}
If each \(\hat{T}_{p_i}(t)\) admits a Hermitian decomposition into Pauli or Clifford basis, and the number of primes \(N(t)\) is poly-logarithmic in system size, then \(\hat{H}_t\) is efficiently simulable on a quantum processor with polynomial depth.

\(\hfill \square\)

